\documentclass[11pt,a4paper]{article}

% -------------------------------------------------------
% Page layout and typography
% -------------------------------------------------------
\usepackage[margin=2.5cm]{geometry}
\usepackage[T1]{fontenc}
\usepackage[utf8]{inputenc}
\usepackage{lmodern}
\usepackage{microtype}
\usepackage{setspace}
\onehalfspacing

% -------------------------------------------------------
% Tables, figures and maths
% -------------------------------------------------------
\usepackage{graphicx}
\usepackage{booktabs}
\usepackage{tabularx}
\usepackage{longtable}
\usepackage{array}
\usepackage{amsmath}
\usepackage{siunitx}

% -------------------------------------------------------
% TikZ figures
% -------------------------------------------------------
\usepackage{tikz}
\usetikzlibrary{
    arrows.meta,
    positioning,
    shapes.geometric,
    shapes.misc,
    calc,
    fit,
    backgrounds
}

% -------------------------------------------------------
% References and links
% -------------------------------------------------------
\usepackage[round,authoryear]{natbib}
\usepackage[hidelinks]{hyperref}
\usepackage[nameinlink,noabbrev]{cleveref}

% -------------------------------------------------------
% Metadata
% -------------------------------------------------------
\title{
Wildfire as an Emerging Systemic Mountain Risk in the European Alps:\\
Processes, Modelling Requirements and Cascading Impacts
}

\author{
Author Name$^{1,*}$ \\
\small $^{1}$Affiliation, Department, Institution, City, Country \\
\small $^{*}$Corresponding author: email@domain.org
}

\date{}

% -------------------------------------------------------
% Begin document
% -------------------------------------------------------
\begin{document}

\maketitle

\subsection{Review objective and research questions}

\begin{enumerate}
    \item xxxxxxxxxxx
    \item yyyyyyyyyyy

\end{enumerate}

\begin{abstract}
xxxxxxxxxxxxxxxxxxxxxxxxxxxxxxxxxxxxxxxxxxxxxxxxxxxxxxxxxxxxxxxxxxxxxxxxxxxxxxxxx 
\end{abstract}

\noindent\textbf{Keywords:}
xxxxxxxxxxxxxxxxxxxxxxxxxxxxxxxxxxxxxxxxxxxxxxxxxxxxxxxxxxxxxxxxxxxxxxxxxxxxxxxxx 

\tableofcontents
\newpage


% =======================================================
\section{Introduction}
% =======================================================

\subsection{Wildfire risk beyond Mediterranean Europe}

\subsection{Why mountain fires are a distinct risk problem}

\subsection{European Alpine regions as an emerging fire-risk frontier}

\subsection{Review objective}


% =======================================================
\section{Review Design and Scope}
% =======================================================

\subsection{Geographical and thematic scope}

\subsection{Literature identification and screening procedure}

\subsection{Thematic synthesis approach}


% =======================================================
\section{The European Alps as a Fire-Prone Mountain System}
% =======================================================

\subsection{Climatic gradients, elevation and seasonal constraints}

\subsection{Land-use legacies, forest expansion and fuel continuity}

\subsection{Differences between Alpine wildfire and lowland or Mediterranean wildfire}

\subsection{Ecosystem, infrastructure and exposed Alpine communities}


% =======================================================
\section{Climate Change and the Reorganisation of Alpine Fire Regimes}
% =======================================================

\subsection{From short-term fire danger to long-term fire-regime change}

\subsection{Warming \& drought}

\subsection{Upslope expansion of fire activity}

\subsection{Changing fire seasonality}

\subsection{Extreme fire-weather synchrony across European mountain regions}


% =======================================================
\section{Mountain Meteorology and Ignition Processes}
% =======================================================

\subsection{Human ignitions in accessible valleys and wildland--urban interfaces}

\subsection{Lightning as a mountain-specific ignition source}

\subsection{Foehn winds, valley winds and slope-wind systems}

\subsection{Wind channelling, turbulence and rapid changes in fire behaviour}


% =======================================================
\section{Dynamic Alpine Fuels and Forest Composition}
% =======================================================

\subsection{Static fuel maps versus dynamic fuel systems}

\subsection{Climate-driven shifts in tree-species suitability}

\subsection{Forest vulnerability, drought stress and disturbance interactions}

\subsection{Land use, treeline dynamics and land abandonment}

\subsection{Understorey, dead biomass and fuel continuity}

\subsection{Implications for future flammability and fire spread}


% =======================================================
\section{Topography, Fire Behaviour and Suppression Constraints}
% =======================================================

\subsection{Slope and Aspect features as drivers of fire behaviour}

\subsection{Elevation, snow cover and seasonal drying windows}

\subsection{Terrain complexity, access constraints and initial attack limits}

\subsection{Weather-only fire danger indices enough in Alpine terrain?}


% % =======================================================
% \section{Soil Dynamics and Post-Fire Geomorphic Cascades}
% % =======================================================

% \subsection{Fire effects on mountain soils}

% \subsection{Soil moisture, organic matter and water repellency}

% \subsection{Runoff generation and post-fire erosion}

% \subsection{Debris flows, shallow landslides and sediment connectivity}

% \subsection{Rockfall, avalanches and loss of protection-forest function}

% \subsection{Time windows of elevated geomorphic susceptibility after fire}


% =======================================================
\section{Modelling Alpine Wildfire Hazard and Susceptibility}
% =======================================================

\subsection{Fire-weather indices and operational danger systems}

\subsection{Ignition models}

\subsection{Spread models in steep and heterogeneous terrain}

\subsection{Climate-scenario hazard models}

\subsection{Dynamic vegetation and fuel-transition models}

\subsection{Uncertainty, transferability and validation in mountain settings}


% =======================================================
\section{From Hazard Mapping to Systemic Risk Assessment}
% =======================================================

\subsection{Limits of single-hazard wildfire assessment}

\subsection{Wildfire as a trigger of cascading Alpine hazards}

\subsection{Compound climate--fire--geomorphic risk}

\subsection{Exposure of settlements, transport corridors, tourism and critical infrastructure}

\subsection{Bias, variance and model complexity in Alpine wildfire risk assessment}

\subsection{Toward dynamic, multi-hazard and uncertainty-aware risk models}


% =======================================================
\section{Governance, Management and Cross-Border Risk}
% =======================================================

\subsection{Fire suppression under rapid spread and difficult access}

\subsection{Prevention, fuel management and fire-smart forestry}

\subsection{Prescribed fire and its limits in Alpine landscapes}

\subsection{Cross-border early warning and emergency coordination}

\subsection{Governance gaps between forest management, civil protection and hazard mitigation}

\subsection{Economic and social implications for Alpine communities}


% % =======================================================
% \section{Synthesis: What Is Known, What Is Missing}
% % =======================================================

% \subsection{Established knowledge}

% \subsection{Emerging but incomplete evidence}

% \subsection{Critical knowledge gaps}

% \subsection{Priority modelling needs}

% \subsection{Priority monitoring needs}

% \subsection{Priority governance needs}


% % =======================================================
% \section{A Conceptual Framework for Alpine Wildfire Systemic Risk}
% % =======================================================

% \subsection{Climate forcing}

% \subsection{Ignition processes}

% \subsection{Fuel and vegetation dynamics}

% \subsection{Topographic controls on fire behaviour}

% \subsection{Post-fire geomorphic cascades}

% \subsection{Exposure, vulnerability and institutional response}

% \subsection{Feedbacks and uncertainty}


% =======================================================
\section{Discussion}
% =======================================================

\subsection{Why Alpine wildfire risk is qualitatively changing}

\subsection{Implications for fire science}

\subsection{Implications for mountain geomorphology}

\subsection{Implications for climate adaptation}

\subsection{Implications for risk governance}


% =======================================================
\section{Conclusions}
% =======================================================

\subsection{Main findings}

\subsection{Scientific contribution}

\subsection{Research agenda}


% =======================================================
\section*{Acknowledgements}
\addcontentsline{toc}{section}{Acknowledgements}

% Text here.


\section*{Funding}
\addcontentsline{toc}{section}{Funding}

% Text here.


\section*{Competing Interests}
\addcontentsline{toc}{section}{Competing Interests}

% Text here.


\section*{Data and Code Availability}
\addcontentsline{toc}{section}{Data and Code Availability}

% Text here.


% =======================================================
\appendix
% =======================================================

\section{PRISMA Search Protocol}

\subsection{Information sources}

\subsection{Search strings}

\subsection{Eligibility criteria}

\subsection{Screening and coding procedure}

\subsection{PRISMA flow diagram}


\section{Data Extraction Template}


\section{Quality Appraisal Criteria}


\section{Glossary of Terms}


% =======================================================
% Bibliography
% =======================================================
\bibliographystyle{apalike}
\bibliography{references}

\end{document}